\chapter{Requirements Analysis}
\label{chap:requirements}

The requirements for the platform are derived from the CZU deployment context. A gap analysis
against the baseline Time.IO system then motivates the scope of the contributions.
Section~\ref{sec:req-usecases} identifies the CZU use cases and stakeholder
needs. Sections~\ref{sec:req-functional} and~\ref{sec:req-nonfunctional} enumerate the functional and non-functional
requirements respectively. Section~\ref{sec:req-gap} compares those requirements against the capabilities of the
original UFZ Time.IO system.

%% ================================================================
\section{CZU Use Cases and Stakeholder Needs}
\label{sec:req-usecases}
%% ================================================================

Czech University of Life Sciences Prague (CZU) operates environmental monitoring infrastructure
across river catchment monitoring projects and experimental plots. The primary stakeholders are
research groups responsible for individual monitoring projects, each of which manages a distinct
set of sensors and requires isolated data views with controlled access. Three actor roles interact
with the platform: \emph{Researcher} (creates projects, configures sensors, analyzes data),
\emph{Administrator} (manages users, groups, and platform-wide settings), and
\emph{External System} (automated data sources such as MQTT devices, SFTP servers, and external
APIs). An additional implicit actor is the \emph{TSM Ingestion Pipeline}, which provisions
infrastructure and processes data on behalf of the other actors.

The use cases are organized into three groups: data ingestion, platform management, and data
consumption. These groups are illustrated in Figures~\ref{fig:uc-ingestion},
\ref{fig:uc-management}, and~\ref{fig:uc-consumption}.

%% PlantUML source files live in figures/*.puml
%% Regenerate PDFs: plantuml -tpdf figures/*.puml

%% --- Use Case Diagram: Data Ingestion ---
\begin{figure}[htbp]
  \centering
  \includegraphics[width=0.85\textwidth]{figures/uc-ingestion.pdf}
  \caption{Use case diagram: data ingestion.}
  \label{fig:uc-ingestion}
\end{figure}

%% --- Use Case Diagram: Platform Management ---
\begin{figure}[htbp]
  \centering
  \includegraphics[width=0.85\textwidth]{figures/uc-management.pdf}
  \caption{Use case diagram: platform management.}
  \label{fig:uc-management}
\end{figure}

%% --- Use Case Diagram: Data Consumption ---
%% NOTE: UC-16 (Run computation) removed — feature decommissioned.
\begin{figure}[htbp]
  \centering
  \includegraphics[width=0.85\textwidth]{figures/uc-consumption.pdf}
  \caption{Use case diagram: data consumption.}
  \label{fig:uc-consumption}
\end{figure}

The identified use cases are summarized in Table~\ref{tab:usecases}.
%% UC-16 (Run computation) removed — feature decommissioned.

\begin{longtable}{l p{0.58\textwidth} p{0.18\textwidth}}
  \caption{Summary of use cases derived from the CZU deployment context.}
  \label{tab:usecases} \\
  \toprule
  ID & Description & Primary actor \\
  \midrule
  \endfirsthead
  \toprule
  ID & Description & Primary actor \\
  \midrule
  \endhead
  \midrule
  \multicolumn{3}{r}{\emph{Continued on next page}} \\
  \endfoot
  \bottomrule
  \endlastfoot
  UC-1  & Upload a CSV data file to a Thing's storage bucket for batch ingestion.
          The file is parsed using the Thing's configured parser and the resulting
          Observations are inserted into the per-thing schema. & Researcher \\
  UC-2  & Stream real-time MQTT data from a field sensor. The device publishes
          payloads to a topic namespace; the mqtt-ingest worker parses
          and stores the Observations. & MQTT Device \\
  UC-3  & Periodically synchronise files from an external SFTP server into a
          Thing's MinIO bucket, triggering file-ingest processing. & SFTP Server \\
  UC-4  & Periodically pull data from an external REST API (e.g.\ DWD, TTN)
          using a built-in or custom syncer script. & External API \\
  UC-5  & Onboard multiple sensors at once by uploading a CSV template containing
          Thing definitions, with per-row validation and error reporting. & Researcher \\
  UC-6  & Create, update, and delete projects. Assign members with roles
          (owner, editor, viewer) and link sensors to projects. & Researcher \\
  UC-7  & Configure a sensor: select ingest type, assign a parser,
          set MQTT credentials, define Datastream metadata, and specify
          geographic coordinates. & Researcher \\
  UC-8  & Define alert rules on a Datastream (threshold, no-data, computation
          failure) with severity levels (info, warning, critical).
          Acknowledge and resolve triggered alerts. & Researcher \\
  UC-9  & Configure QA/QC test pipelines per project using SaQC functions with
          tunable parameters. Trigger evaluation on demand. & Researcher \\
  UC-10 & Create and manage Keycloak groups and assign users to role-based
          subgroups (viewers, editors, admins). & Administrator \\
  UC-11 & Create and publish geospatial layers (vector, raster, WMS) through
          GeoServer with configurable access control. & Administrator \\
  UC-12 & Visualize time-series data for one or more Datastreams with
          interactive charts, decimation controls, and adjustable Y-axis bounds.
          & Researcher \\
  UC-13 & View sensor locations on an interactive map with color-coded markers
          and overlay geospatial layers. & Researcher \\
  UC-14 & Export Observation data as CSV for offline analysis. & Researcher \\
  UC-15 & Query the OGC SensorThings API endpoint for machine-to-machine data
          access using standard \texttt{\$filter} and \texttt{\$expand} queries. & External Client \\
\end{longtable}

%% ================================================================
\section{Functional Requirements}
\label{sec:req-functional}
%% ================================================================

The functional requirements are derived from the use cases in Section~\ref{sec:req-usecases}
and are grouped into eight areas. Each requirement is assigned an identifier for traceability
in Chapter~\ref{chap:testing}.

\subsection{Data Ingestion}
\label{sec:req-functional-ingestion}

\begin{description}
  \item[FR-1] The system shall ingest real-time sensor data via MQTT, supporting device-specific
    parsers that can be registered and selected per Thing. \emph{(UC-2)}
  \item[FR-2] The system shall ingest historical data from CSV files uploaded to a Thing's
    object-storage bucket, parsed using the Thing's configured file parser. \emph{(UC-1)}
  \item[FR-3] The system shall periodically synchronise data from external SFTP servers,
    downloading new files into the Thing's bucket for subsequent file-ingest processing.
    \emph{(UC-3)}
  \item[FR-4] The system shall periodically pull data from external REST APIs using built-in
    syncers (DWD, TTN, Bosch) or user-provided custom syncer scripts. \emph{(UC-4)}
  \item[FR-5] The system shall support uploading custom MQTT parser scripts and custom API
    syncer scripts, validated before deployment. \emph{(UC-2, UC-4)}
\end{description}

\subsection{Sensor Management}
\label{sec:req-functional-sms}

\begin{description}
  \item[FR-6] The system shall support creating, updating, and deleting Things with configurable
    ingest type (MQTT, SFTP, external API), parser assignment, and geographic coordinates.
    \emph{(UC-7)}
  \item[FR-7] The system shall allow bulk import of Thing definitions from a CSV file, with
    per-row validation and error reporting. A downloadable CSV template shall be provided.
    \emph{(UC-5)}
  \item[FR-8] The system shall track sensor activity per Thing, detecting inactivity based on
    a configurable threshold (default: 24~hours) and raising alerts when no data is received.
    \emph{(UC-8)}
  \item[FR-9] The system shall manage MQTT device types and CSV parser definitions as reusable
    catalog entries that can be assigned to multiple Things. \emph{(UC-7)}
\end{description}

\subsection{Project Management}
\label{sec:req-functional-projects}

\begin{description}
  \item[FR-10] The system shall support creating, updating, and deleting projects. Each project
    shall group a set of Things with shared access control. \emph{(UC-6)}
  \item[FR-11] The system shall support assigning users to projects with one of three roles:
    owner (full control), editor (modify data and configuration), or viewer (read-only access).
    \emph{(UC-6)}
  \item[FR-12] The system shall create a dedicated TimeIO database schema per project for
    per-thing observation storage. \emph{(UC-6)}
\end{description}

\subsection{Data Quality}
\label{sec:req-functional-quality}

\begin{description}
  \item[FR-13] The system shall support configuring QA/QC test pipelines per project, using
    functions from the SaQC (Statistical Quality Control) library with tunable parameters.
    \emph{(UC-9)}
  \item[FR-14] The system shall execute QA/QC tests automatically when new data is parsed
    (triggered by the \texttt{data\_parsed} MQTT event) and store quality flags alongside
    Observations. \emph{(UC-9)}
  \item[FR-15] The system shall allow on-demand triggering of QA/QC evaluation through the
    API. \emph{(UC-9)}
\end{description}

\subsection{Visualization}
\label{sec:req-functional-viz}

\begin{description}
  \item[FR-16] The system shall provide interactive time-series charts with support for multiple
    Datastream overlays, decimation controls for large datasets, and adjustable Y-axis bounds.
    \emph{(UC-12)}
  \item[FR-17] The system shall display sensor locations on an interactive map with color-coded
    markers and popup details. WMS/WFS geospatial layers shall be overlayable. \emph{(UC-13)}
  %% FR-18 removed — custom dashboards decommissioned.
  \item[FR-19] The system shall allow exporting Observation data as CSV files for offline
    analysis. \emph{(UC-14)}
\end{description}

\subsection{Authentication and Access Control}
\label{sec:req-functional-auth}

\begin{description}
  \item[FR-20] The system shall authenticate users via OpenID Connect (OIDC) through a
    Keycloak identity provider. User registration, login, token refresh, and password change
    shall be supported. \emph{(UC-6, UC-10)}
  \item[FR-21] The system shall enforce a two-tier role-based access control model:
    Tier~1 at the group level (viewer, editor, admin roles stored in Keycloak subgroups) and
    Tier~2 at the project level (owner, editor, viewer roles stored in the application database).
    Realm-level administrators shall override all restrictions. \emph{(UC-10)}
  \item[FR-22] The system shall support creating and managing Keycloak groups with role-based
    subgroups. Group membership shall propagate as default permissions for projects within
    the group. \emph{(UC-10)}
\end{description}

\subsection{Batch Import}
\label{sec:req-functional-batch}

\begin{description}
  \item[FR-23] The system shall allow uploading historical time-series data as CSV files
    (up to 200~MB) with automatic Datastream mapping. \emph{(UC-1)}
  \item[FR-24] The system shall allow importing geospatial data in GeoJSON format (up to
    200~MB) with streamed processing. \emph{(UC-11)}
\end{description}

\subsection{Multilingual Interface}
\label{sec:req-functional-i18n}

\begin{description}
  \item[FR-25] The user interface shall be available in English, Czech, and Slovak, with
    locale detection based on browser settings and manual override. \emph{(all use cases)}
\end{description}

%% ================================================================
\section{Non-Functional Requirements}
\label{sec:req-nonfunctional}
%% ================================================================

The non-functional requirements govern the operational and quality properties of the platform.

\begin{description}
  \item[NFR-1] \textbf{Containerized deployment.} The entire platform shall be deployable as
    Docker Compose services. Rootless Podman shall be supported as an alternative container
    runtime.

  \item[NFR-2] \textbf{OGC SensorThings API compliance.} The platform shall expose observation
    data through a FROST-Server instance that is fully compliant with OGC STA Part~1 (Sensing),
    including \texttt{\$filter}, \texttt{\$expand}, and \texttt{\$orderby} query parameters.

  \item[NFR-3] \textbf{OIDC standards compliance.} Authentication shall follow the OpenID Connect
    standard with JWT token validation, JWKS key rotation, and standard token endpoints.

  \item[NFR-4] \textbf{Modular extensibility.} New ingestion worker types, parsers, and API
    syncers shall be addable without modifying the core platform. Custom parsers and syncers shall
    be uploadable as Python scripts at runtime.

  \item[NFR-5] \textbf{Security.} User-uploaded parser and syncer scripts shall be sandboxed
    through AST-based validation that blocks dangerous operations (\texttt{eval}, \texttt{exec},
    \texttt{open}, \texttt{getattr}/\texttt{setattr}). API rate limiting shall protect
    authentication endpoints. Credentials stored in the database shall be encrypted.

  \item[NFR-6] \textbf{Maintainability.} The codebase shall use database migrations (Alembic
    for \texttt{water-dp}, Flyway for TSM) to manage schema evolution. The API shall generate
    OpenAPI documentation automatically.

  \item[NFR-7] \textbf{eLTER interoperability alignment.} The architecture shall support
    future extension toward eLTER metadata standards and cross-site federation without
    requiring a fundamental redesign.

  \item[NFR-8] \textbf{Unified deployment.} The combined TSM and \texttt{water-dp} stacks shall
    be deployable through a single configuration file and Makefile, with automated secret
    generation and health checks.
\end{description}

%% ================================================================
\section{Gap Analysis: Time.IO vs.\ Target Platform}
\label{sec:req-gap}
%% ================================================================

Table~\ref{tab:gap-analysis} summarizes the gap between the Time.IO baseline capabilities and the
requirements identified in Sections~\ref{sec:req-functional} and~\ref{sec:req-nonfunctional}.
Each row identifies a capability that the upstream Time.IO system does not provide, the specific
requirements it blocks, and the component that addresses the gap in this thesis.

\begin{longtable}{l p{0.58\textwidth} p{0.19\textwidth}}
  \caption{Gap analysis: Time.IO baseline vs.\ target platform requirements.}
  \label{tab:gap-analysis} \\
  \toprule
  Gap & Description & Addressed by \\
  \midrule
  \endfirsthead
  \toprule
  Gap & Description & Addressed by \\
  \midrule
  \endhead
  \midrule
  \multicolumn{3}{r}{\emph{Continued on next page}} \\
  \endfoot
  \bottomrule
  \endlastfoot

  G-1 & \textbf{No project-oriented data model.} Time.IO organizes Things by
        configdb projects (Keycloak groups), but provides no application-level
        project entity with member roles or sensor assignment.
        \emph{(FR-10--12)} & \texttt{water-dp} \\[6pt]

  G-2 & \textbf{No user-facing web portal.} The upstream frontends
        (\texttt{tsm-frontend}, \texttt{thing-management}) provide only basic
        Thing CRUD through Django or Vue.js admin interfaces. No unified portal
        with time-series visualization, map display, or alerting exists.
        \emph{(FR-16--19)} & \texttt{water-dp} \\[6pt]

  G-3 & \textbf{No geospatial layer management.} Time.IO stores geographic
        coordinates in the configdb but provides no GIS integration, layer
        publishing, or spatial query capabilities.
        \emph{(FR-17, FR-24)} & \texttt{water-dp} \\[6pt]

  G-4 & \textbf{Non-functional local STA views.} The original local STA views
        contained issues that prevented FROST-Server from returning valid
        responses for certain entity types.
        \emph{(NFR-2)} & TSM adaptation \\[6pt]

  G-5 & \textbf{Non-functional SMS service.} The Sensor Management System
        backend was not operational; the \texttt{timeio-and-sms} combined
        deployment was impossible to start due to broken dependencies.
        \emph{(FR-6, FR-9)} & TSM adaptation \\[6pt]

  G-6 & \textbf{No alert evaluation subsystem.} Time.IO provides no mechanism
        for threshold monitoring, no-data detection, or notification of
        anomalous conditions.
        \emph{(FR-8, UC-8)} & \texttt{water-dp} \\[6pt]

  G-7 & \textbf{No unified deployment mechanism.} The TSM and application
        stacks are separate Docker Compose deployments with no shared
        configuration, secret management, or health-check tooling.
        \emph{(NFR-1, NFR-8)} & \texttt{hydro-deploy} \\[6pt]

  G-8 & \textbf{No Podman rootless support.} The upstream Compose files use
        Docker-specific features (environment list format, disabled service
        references) that are incompatible with rootless Podman.
        \emph{(NFR-1)} & \texttt{hydro-deploy} \\[6pt]

  G-9 & \textbf{No multilingual interface.} The upstream frontends are
        English-only with no internationalization framework. Czech and Slovak
        language support is required for CZU users.
        \emph{(FR-25)} & \texttt{water-dp} \\[6pt]

  %% G-10 removed — computation framework decommissioned.

  G-11 & \textbf{Upstream frontend bugs.} Both \texttt{tsm-frontend} (Django)
         and \texttt{thing-management} (Vue.js) generated MQTT configuration
         messages that did not match the ingestion pipeline's expectations,
         causing silent provisioning failures.
         \emph{(FR-6, FR-7)} & \texttt{water-dp} \\[6pt]

  G-12 & \textbf{Hardcoded parsers and device types.} All file parsers, MQTT
         device-type parsers, and API syncers are implemented as Python classes
         compiled into the Docker image. Adding support for a new sensor model
         or file format requires source-code changes and redeployment; there is
         no runtime mechanism for users to register custom parsers or device
         types through a management interface.
         \emph{(FR-5, FR-9, NFR-4)} & \texttt{water-dp} \\

\end{longtable}

The eleven gaps identified in Table~\ref{tab:gap-analysis} define the scope of the three
contributions described in Chapters~\ref{chap:architecture} and~\ref{chap:implementation}:
the \texttt{water-dp} application layer addresses gaps~G-1--G-3 and G-6--G-9, G-11, and G-12,
the TSM adaptations address gaps~G-4 and~G-5, and the \texttt{hydro-deploy} deployment
orchestrator addresses gaps~G-7 and~G-8. Each gap is revisited in the requirements evaluation
in Chapter~\ref{chap:testing}.
